\begin{table}[htbp]
\centering
\caption{Escore médio dos componentes por grupo e teste de Kruskal-Wallis entre grupos}
\label{tab:escores}
\resizebox{\textwidth}{!}{%
\begin{tabular}{l r r r c c c}
\hline
\textbf{Componente} & \textbf{A} & \textbf{B} & \textbf{C} & $H$ & $p$ & $\varepsilon^2$ \\ \hline
D1 PC1 (bens duráveis) & 1,67 & -1,51 & 0,60 & 31,1 & $<0{,}001$ & 0,63 \\
D2 PC1 (formalização e escolaridade) & 1,18 & -0,71 & 0,00 & 13,4 & 0,001 & 0,25 \\
D3 PC1 (infraestrutura pública) & 2,68 & -1,33 & -0,33 & 25,5 & $<0{,}001$ & 0,51 \\
D3 PC2 (condição da moradia) & -0,46 & -1,03 & 1,54 & 23,0 & $<0{,}001$ & 0,46 \\
D4 PC1 (conectividade) & 1,70 & -1,38 & 0,43 & 30,0 & $<0{,}001$ & 0,61 \\
\hline
\end{tabular}}
\par\footnotesize Fonte: elaboração própria. Média dos escores dos segmentos de cada grupo (grupos A a C). Os escores têm média zero e desvio-padrão de 1,2 a 2,2 entre os 49 segmentos, e o sinal de cada componente é arbitrário (a leitura segue as cargas da Seção~\ref{sec:res_pca}). $H$: estatística de Kruskal-Wallis (4 graus de liberdade); $\varepsilon^2$: tamanho de efeito. Os $p$-valores são descritivos, porque os grupos foram formados com os mesmos componentes (Seção~\ref{sec:metodo_robustez}); todos permanecem abaixo de 0,01 após a correção de Benjamini-Hochberg.
\end{table}
